\section{Introduction}

\begin{frame}%[label=wip]
    \titlepage
    \note{%
        \textbf{Before you start}
        \begin{itemize}
            \item Check that everyone got the template link.
            \item Nobody has to keep up.
                Every block ends with a checkpoint file that puts them back on track.
        \end{itemize}
    }
\end{frame}

\begin{frame}
    \frametitle{In two hours, you will have (hopefully) written this}
    \begin{center}
        \includegraphics[page=1, height=.5\textheight, frame]{renders/checkpoint-04.pdf}
        \hfill
        \includegraphics[page=2, height=.5\textheight, frame]{renders/checkpoint-04.pdf}
        \hfill
        \includegraphics[page=3, height=.5\textheight, frame]{renders/checkpoint-04.pdf}
        \hfill
        \includegraphics[page=4, height=.5\textheight, frame]{renders/checkpoint-04.pdf}

        \medskip
        {\small\color{black!60} Title, table of contents, equations, a table, a figure, a bibliography.}
    \end{center}
    \note{%
        \begin{itemize}
            \item Short document containing basic LaTeX stuff.
        \end{itemize}
    }
\end{frame}

\begin{frame}
    \frametitle{Who is in the Room?}
    \vfill
    \begin{tcbraster}[%
        raster columns=2,
        raster equal height,
        raster column skip=8mm,
        raster row skip=5mm,
        deckbox,
        thinframe,
        colback=white,
        fontupper=\large,
    ]
        \begin{tcolorbox}
            {\Huge 1}\par I have never heard of \LaTeX{}
        \end{tcolorbox}
        \begin{tcolorbox}
            {\Huge 2}\par I have heard of it, never used it
        \end{tcolorbox}
        \begin{tcolorbox}
            {\Huge 3}\par I have compiled something once
        \end{tcolorbox}
        \begin{tcolorbox}
            {\Huge 4}\par I use it already
        \end{tcolorbox}
    \end{tcbraster}
    \vfill
    \note{%
        Hands up, 30 seconds.
        Last year roughly 95\,\% were in groups 1 and 2, so the deck is built for them.
        \begin{itemize}
            \item Say that student should tell me if I need to slow down.
        \end{itemize}
    }
\end{frame}

\begin{frame}
    \frametitle{Sound Familiar?}
    \vfill
    \large
    \begin{itemize}
        \item<1-> New heading style, but only on the heading you selected
        \item<2-> One image moved, the next three pages look different
        \item<3-> Captions that all look a bit different
    \end{itemize}
    \vfill
    \note{%
        \begin{itemize}
            \item Ask: who has written a long document in Word or Google Docs?
                Who had one of these problems?
            \item None of this is Word's fault alone.
                Word has styles, but it also lets you format text directly,
                so a consistent result depends on your discipline.
        \end{itemize}
    }
\end{frame}

% + and - in front of the items, coloured like the rest of the deck.
\newcommand{\pro}{\textcolor{structure.fg}{\textbf{+}}}
\newcommand{\con}{\textcolor{structure.fg}{\textbf{--}}}

\begin{frame}
    \frametitle{Two Ways to Write a Document}
    \begin{columns}[T, onlytextwidth]
        \begin{column}{.47\textwidth}
            \begin{block}{WYSIWYG}
                {\small\color{black!60} Word, Google Docs, \ldots\par}
                \smallskip
                \begin{itemize}
                    \item[\pro] You see the result while you type
                    \item[\pro] Easy to start
                    \item[\con] You keep the formatting consistent
                    \item[\con] Large documents get hard to manage
                \end{itemize}
            \end{block}
        \end{column}
        \begin{column}{.47\textwidth}
            \uncover<2->{%
                \begin{block}{\LaTeX{}}
                    {\small\color{black!60} Text and commands in a plain-text file\par}
                    \smallskip
                    \begin{itemize}
                        \item[\pro] Formatting stays consistent
                        \item[\pro] Numbers and references update automatically
                        \item[\con] You compile to see the result
                        \item[\con] Takes time to learn
                    \end{itemize}
                \end{block}
            }
        \end{column}
    \end{columns}
    \note{%
        \begin{itemize}
            \item WYSIWYG: ``what you see is what you get''.
            \item Be fair to Word.
                It has styles, automatic captions and an equation editor.
                The difference: Word \emph{lets} you format directly, \LaTeX{} makes you describe the structure.
            \item More points (Git, math, collaboration) are in the handout.
        \end{itemize}
    }
\end{frame}


\begin{frame}
    \frametitle{When it is Worth the Trouble}
    \begin{center}
        % =====================================================================
%  Effort vs. complexity of a document, for the slides and the handout.
%  A sketch, not measured data.
%  Needs: pgfplots, hologo
%  Colours: `figaccent' (LaTeX) and `figalt' (word processor), the deck's
%  colours unless defined before \input-ing this file.
% =====================================================================
\providecolor{figaccent}{HTML}{8C195F}
\providecolor{figalt}{HTML}{746BD2}
\begin{tikzpicture}
    \begin{axis}[
        smooth, samples=100,
        axis lines=left,
        axis line style={thick},
        width=.85\textwidth,
        height=4.6cm,
        xlabel=Complexity of the document,
        ylabel=Effort,
        x label style={at={(axis description cs:1,.1)}, anchor=east},
        ytick=\empty,
        xtick={2, 4, 6, 11, 13, 16, 20, 22},
        xticklabels={Text, Figures, Tables, Maths, Plots, {Technical\\ Report}, Thesis, Book},
        xticklabel style={align=center},
        x tick label style={rotate=70, anchor=east, font=\footnotesize},
        legend style={legend pos=north west},
        no markers, domain=0:25,
    ]
        \addplot[densely dashed, very thick, figalt, domain=0:14]{exp(x/3 - 2.5)};
        \addlegendentry{Word processor}
        \addplot[very thick, figaccent]{3/(1 + exp(-2*(x-2.5))) + x/50*exp((x-4)/12)};
        \addlegendentry{\textrm{\hologo{LaTeX}}}
    \end{axis}
\end{tikzpicture}

    \end{center}
    \begin{columns}[t, onlytextwidth]
        \begin{column}{.47\textwidth}
            \begin{alertblock}<2->{}
                One-page note: word processor
            \end{alertblock}
        \end{column}
        \begin{column}{.47\textwidth}
            \begin{exampleblock}<2->{}
                Report, thesis: \LaTeX{}
            \end{exampleblock}
        \end{column}
    \end{columns}
    \note{%
        \begin{itemize}
            \item Take this plot with a grain of salt.
        \end{itemize}
    }
\end{frame}

\begin{frame}
    \frametitle{How \LaTeX{} Works}
    \vfill
    \begin{center}
        \begin{tikzpicture}[
            >=stealth,
            stage/.style={
                draw=structure.fg, line width=1pt, rounded corners=1.3mm,
                fill=block body example.bg,
                minimum width=3.2cm, minimum height=1.6cm, align=center,
            },
            flow/.style={->, line width=1.2pt, structure.fg},
        ]
            \node[stage] (src) at (0,0)   {Source file\\[1mm]\texttt{report.tex}};
            \node[stage] (tex) at (4.6,0) {\LaTeX{}};
            \node[stage] (pdf) at (9.2,0) {PDF\\[1mm]\texttt{report.pdf}};
            \node[above, structure.fg] at (tex.north) {compile};
            \draw[flow] (src) -- (tex);
            \draw[flow] (tex) -- (pdf);
            \draw[flow] (pdf.south) to[out=-120, in=-60]
                node[below] {change the source, compile again} (src.south);
        \end{tikzpicture}
    \end{center}
    \vfill
    \note{%
        \begin{itemize}
            \item The source file is plain text.
                Any editor can open it.
            \item You never edit the PDF.
                You change the source and compile again.
            \item The program that compiles is called the engine,
                e.g.\ \hologo{pdfLaTeX} or \hologo{LuaLaTeX}.
                There are more, but out of scope for now.
        \end{itemize}
    }
\end{frame}

% Overlay 1: article, overlay 2: scrartcl. The renders are the listings
% compiled as they are; the viewport shows the same region of the page
% for both (in bp from the bottom-left corner of the A4 page).
\begin{frame}
    \frametitle{Same Text, One Word Changed}
    \only<1| handout:1>{\tcbinputlst[fontsize=\footnotesize]{listings/00-class-article.tex}}%
    \only<2| handout:2>{\tcbinputlst[fontsize=\footnotesize, highlightlines={1}]{listings/00-class-scrartcl.tex}}%
    \begin{tcolorbox}[deckbox, thinframe, colback=white, left=0mm, right=0mm, top=0mm, bottom=0mm]
        \only<1| handout:1>{\includegraphics[width=\linewidth, viewport=80 677 470 764, clip]{renders/00-class-article.pdf}}%
        \only<2| handout:2>{\includegraphics[width=\linewidth, viewport=80 677 470 764, clip]{renders/00-class-scrartcl.pdf}}%
    \end{tcolorbox}
    \note{%
        \begin{itemize}
            \item Click back and forth between the two states.
            \item The text did not change, only the document class.
                Heading font, margins and text size follow from it.
            \item Do not explain \texttt{\textbackslash documentclass} here.
                It is the first thing in block 1.
        \end{itemize}
    }
\end{frame}

% ---------------------------------------------------------------------
%  History
% ---------------------------------------------------------------------
% The timeline is shared with the handout: figures/tex-timeline.tex
\begin{frame}
    \frametitle{Not a New Idea}
    \vfill
    \begin{center}
        % =====================================================================
%  Timeline of TeX, LaTeX and the engines, for the slides and the handout.
%  Needs: tikz, hologo
%  Above the line: TeX and the macro packages written for it.
%  Below the line: newer programs (engines) that run them.
%
%  \timelineunit: width of one year. Default 0.4cm (slides);
%  define it before \input-ing this file to change it, e.g. in the handout:
%      \newcommand{\timelineunit}{0.36cm}
%  Colour: `figaccent', the deck's purple unless defined before.
% =====================================================================
\providecommand{\timelineunit}{0.4cm}
\providecolor{figaccent}{HTML}{8C195F}
\begin{tikzpicture}[
    x=\timelineunit,
    tl year/.style={font=\scriptsize\bfseries, text=figaccent},
    tl name/.style={font=\small, align=center},
    tl who/.style={font=\scriptsize, text=black!60, align=center},
]
    % 1975 at x=0, one unit per year.
    % \lvl: 1 or 2 above the line, -1 or -2 below; 2 sits further out,
    % for entries that would otherwise touch their neighbour.
    \draw[line width=1pt, figaccent, ->] (1,0) -- (35,0);
    \foreach \y/\lbl/\person/\lvl in {
            1978/\hologo{TeX}/Donald Knuth/1,
            1984/\hologo{LaTeX}/Leslie Lamport/1,
            1994/\hologo{LaTeXe}/\hologo{LaTeX}\\Project/1,
            1996/\hologo{ConTeXt}/Hans Hagen/2,
            1996/\hologo{pdfTeX}/H\`an Th\'{\^e} Th\`anh/-1,
            2004/\hologo{XeTeX}/Jonathan Kew/-1,
            2007/\hologo{LuaTeX}/\hologo{LuaTeX} team/-2} {
        \pgfmathsetmacro{\x}{\y-1975}
        \pgfmathsetmacro{\out}{(abs(\lvl)-1)*1.25}
        \draw[figaccent, line width=1pt] (\x,-.15) -- (\x,.15);
        % the year always sits under the line; same year twice is fine
        \node[tl year] at (\x,-.4) {\y};
        \ifnum\lvl>0
            \draw[black!30] (\x,.15) -- (\x,.3+\out);
            \node[tl who, anchor=south] (tl-who) at (\x,.3+\out) {\person};
            \node[tl name, anchor=south, inner ysep=1pt] at (tl-who.north) {\lbl};
        \else
            \draw[black!30] (\x,-.65) -- (\x,-.8-\out);
            \node[tl name, anchor=north, inner ysep=1pt] (tl-name) at (\x,-.8-\out) {\lbl};
            \node[tl who, anchor=north] at (tl-name.south) {\person};
        \fi
    }
\end{tikzpicture}

    \end{center}
    \vfill
    \note{%
        \begin{itemize}
            \item Only background: \LaTeX{} is older than most of the room, and still maintained.
                The \LaTeX{} kernel gets new releases twice a year.
            \item Top: \TeX{} and what was written on top of it.
                Bottom: newer programs that run it.
                \hologo{LaTeXe} from 1994 is the \LaTeX{} everyone uses today.
            \item \hologo{ConTeXt} is a separate system with the same roots, not a part of \LaTeX{}.
        \end{itemize}
    }
\end{frame}



\begin{frame}[fragile]
    \frametitle{Getting Started}
    \begin{columns}[t, onlytextwidth]
        \begin{column}{.45\textwidth}
            {\Large\bfseries Today: Overleaf\par}
            \smallskip
            In the browser, nothing to install
            \begin{enumerate}
                \item Account at \textbf{overleaf.com}
                \item New Project, then Upload Project
                \item Drop in the template zip
                \item Open \texttt{workshop.tex}, hit \textbf{Recompile}
            \end{enumerate}
        \end{column}
        \begin{column}{.5\textwidth}
            {\Large\bfseries Later: on your own machine\par}
            \medskip
            \begin{tabular}{@{}lll@{}}
                \toprule
                & \textbf{Install} & \textbf{Edit in} \\
                \midrule
                Windows & \hologo{MiKTeX} & VS Code, TeXstudio \\
                macOS   & Mac\TeX{}       & TeXShop, VS Code \\
                Linux   & \hologo{TeXLive}& VS Code, TeXstudio \\
                \bottomrule
            \end{tabular}

            \bigskip
            Compile with:
            \begin{tcolorbox}[codebg]
                \mintinline{bash}{latexmk -pdf workshop.tex}
            \end{tcolorbox}
        \end{column}
    \end{columns}
    \note{%
        Two minutes, hard limit.
        Do not demo an installation.
        \begin{itemize}
            \item Overleaf today means the same setup for everyone,
                so every problem today is a \LaTeX{} problem.
            \item This slide exists so that nobody spends the workshop installing software,
                and so that people who want a local setup have the names written down.
            \item \texttt{latexmk} matters more than it looks:
                it is why they never have to count compilation passes.
        \end{itemize}
    }
\end{frame}

% The roadmap, revealed one block per click (state 0, see preamble.tex).
\begin{frame}
    \frametitle{Schedule}
    \vfill
    \begin{roadmap}{0}
        \roadmapresults     % comment out to reveal the block names only
    \end{roadmap}
    \vfill
    \note{%
        \begin{itemize}
            \item Walk along the line: four blocks, a break between each.
            \item Say it: ``you will see this slide five more times.''
            \item Timing from here: block 1 ends at about 26 minutes on the clock.
        \end{itemize}
    }
\end{frame}
